\documentclass[11pt]{article}
\usepackage[margin=1in]{geometry}
\usepackage{array}
\usepackage{booktabs}
\usepackage{amsmath,amssymb}
\usepackage{graphicx}
\usepackage[hidelinks]{hyperref}
\newcommand{\code}[1]{\texttt{#1}}

\title{SAC vs.\ SAC+JEPA at a 10{,}000-transition budget}
\author{}
\date{September 20, 2026}

\begin{document}
\maketitle

\noindent
Both runs share the standardized centered-Gaussian terrain, periodic tiling,
random start positions, the $P_{\max}=6$ longitudinal power cap, 16 parallel
environments, the same replay configuration, the same evaluation interval, and
the held-out terrain seeds $100001$--$100004$. Only the actor objective differs.
All numbers below are single-run, 10k-transition, RTX~3090 results.

\section{What differs}

\begin{center}
\begin{tabular}{>{\raggedright\arraybackslash}p{0.24\linewidth}
                >{\raggedright\arraybackslash}p{0.33\linewidth}
                >{\raggedright\arraybackslash}p{0.33\linewidth}}
\toprule
\textbf{Aspect} & \textbf{SAC (vanilla)} & \textbf{SAC+JEPA} \\
\midrule
Actor objective & $\mathbb E[\alpha\log\pi(a|o)-\min(Q_1,Q_2)]$ &
$\mathbb E[\alpha\log\pi(a|o)-\min(Q_1,Q_2)] + \lambda_{\rm JEPA}\mathcal L_{\rm JEPA} + c_{\rm var}\mathcal L_{\rm var}$ \\
Predictive term & none &
one-step latent MSE: $\lVert q_\psi([f_\theta(o_t),a_t])-\operatorname{sg}[\bar f(o_{t+1})]\rVert_2^2$ \\
$\lambda_{\rm JEPA}$ & $0$ & $0.1$ (measured run); $0.025$ (corrected overlay) \\
EMA target cadence & not applicable &
once per replay-gradient update (measured run, unguarded);
once per $64$ collected transitions (corrected) \\
Variance floor & none &
none (measured run); $0.01\,\mathrm{relu}(0.1-\mathrm{std}(f_\theta))$ (corrected) \\
Reported metrics & \code{actor\_loss}, \code{critic\_loss}, \code{alpha} &
adds \code{actor\_total\_loss}, \code{jepa\_loss}, and (corrected) \code{jepa\_variance\_loss}, \code{jepa\_embedding\_std} \\
Replay, collection, terrain, seeds & shared & shared \\
\bottomrule
\end{tabular}
\end{center}

\section{Result}

\begin{center}
\begin{tabular}{lcc}
\toprule
\textbf{Transitions} & \textbf{SAC} & \textbf{SAC+JEPA} \\
\midrule
2{,}048  & $0.064 \pm 0.031$ & $0.064 \pm 0.031$ \\
4{,}096  & $0.190 \pm 0.027$ & $0.190 \pm 0.029$ \\
6{,}144  & $0.292 \pm 0.042$ & $0.279 \pm 0.041$ \\
8{,}192  & $0.673 \pm 0.041$ & $0.641 \pm 0.048$ \\
10{,}000 & $\mathbf{1.655 \pm 0.056}$ & $\mathbf{1.278 \pm 0.042}$ \\
\bottomrule
\end{tabular}
\end{center}

\noindent
The endpoints differ by $0.377$ ($22.8\%$) in favour of vanilla SAC. The
$\pm$ values are standard deviations across the four held-out terrain seeds,
\emph{not} across training runs, so they are not an error bar on the
SAC-vs-SAC+JEPA difference; with one training seed per arm this gap is not
statistically distinguishable.

\begin{figure}[h]
\centering
\includegraphics[width=0.48\linewidth]{../../../dumps/sac_10k/plots/learning_curve.png}
\includegraphics[width=0.48\linewidth]{../../../dumps/sac_jepa_10k_v2/plots/learning_curve.png}
\caption{Left: SAC-only. Right: SAC+JEPA ($\lambda_{\rm JEPA}=0.1$,
pre-correction branch). Both are at the same 10k budget.}
\end{figure}

\section{Why SAC+JEPA is below vanilla}

\begin{enumerate}
\item \textbf{The measured run used the pre-correction EMA cadence.} The target
was updated once per replay-gradient update. With 16 environments and one
update per transition, roughly $1{,}056$ EMA calls occur between transitions
$1{,}000$ and $2{,}048$, so $1-(1-0.01)^{1056}\approx 99.998\%$ of the online
weights are absorbed: the target is numerically the online encoder. The
objective then reduces to a near-identity map, which is exactly what the
observed loss of $9.1\times10^{-6}$ to $3.5\times10^{-5}$ indicates.
\item \textbf{One-step horizon.} With $\mathrm{d}t=0.02$ and two substeps, a
step advances $0.04$\,s, so consecutive field-of-view observations are nearly
identical. Even with a correctly lagged target, ``predict $z_{t+1}$ from
$z_t$'' is close to trivial and carries little dynamics signal.
\item \textbf{The coefficient was inherited from PPO.} PPO+JEPA reached a loss
of $5.95\times10^{-4}$ at the same $\lambda$, about $17\times$ the SAC loss
scale, so $\lambda=0.1$ is not matched to the SAC encoder.
\item \textbf{Placement.} JEPA shapes only the actor encoder. The twin $Q$
critics keep independent encoders and receive no auxiliary signal, so the term
cannot improve value estimation and can only perturb the policy
representation.
\item \textbf{Budget and variance.} At $10$k transitions both curves are inside
their steep takeoff (SAC rises $0.67\rightarrow1.65$ between $8{,}192$ and
$10{,}000$), so a one-step shift in takeoff timing moves the endpoint by
roughly the observed margin.
\end{enumerate}

\section{What the corrected branch changes, and what remains}

The corrected code (i) updates the EMA target once per $64$ collected
transitions, decoupling the lag from \code{num\_parallel\_envs} and
\code{updates\_per\_step} (effective lag $\approx 637$ gradient steps instead
of $\approx 100$); (ii) adds a VICReg-style variance floor on the online
representation; (iii) logs the unregularized actor objective separately from
the optimized combined objective; and (iv) begins at $\lambda_{\rm JEPA}=0.025$.

Two limitations remain. First, \textbf{the JEPA term is negligible relative to
the variance term} at the corrected scales: at the observed 10k loss of
$9.1\times10^{-6}$ the JEPA contribution is $2.3\times10^{-7}$ while the
variance term is $\approx 5\times10^{-4}$, a ratio of roughly $2{,}200$ (and
still $\approx 570$ at the largest observed loss). A corrected run therefore
tests a variance regularizer more than it tests JEPA. Second, the trainer
\emph{raises} when $\lambda_{\rm JEPA}>0$ and $c_{\rm var}\le 0$, so the
variance term cannot be ablated, and the same guard is absent from the PPO
trainer. The horizon is still one step and the objective is still attached only
to the actor encoder.

\medskip
\noindent
\textbf{Status.} No 10k run of the corrected branch exists yet. The $1.278$
result above came from the earlier branch, so it is evidence that the original
auxiliary path ran, not a test of the corrected objective. Multi-seed runs at
$\ge 100$k transitions, with the variance term ablated, are required before any
performance conclusion.

\end{document}
